\section{Our MILP Model}
In this study, we adopt the related-key differential cryptanalysis model introduced in \cite{halil2026milp}. To facilitate a comprehensive comparison, 8-, 10-, and 12-round instances of the model are employed to evaluate the performance of various solvers. For the sake of completeness, a brief overview of related-key differential cryptanalysis and the ITUbee cipher is also provided; for a more detailed exposition, the reader is referred to \cite{halil2026milp}.

\subsection{ITUbee}

ITUbee \cite{karakoc2013itubee} is a lightweight software-oriented block cipher specifically designed for resource-constrained devices which include a microcontroller and have a limited battery power such as sensor nodes in wireless sensor networks.

The cipher has Fiestel structure which consists of 20 rounds. İt has 80 bits block size and 80 bits key size. The round structure of the cipher is given in Figure \ref{fig:itubee}. Cipher also has a key whitening layer in the beginning and end of the 20 rounds. The encryption process of the cipher is given in Algorithm \ref{alg:ITUbee}. 

\begin{figure}[!ht]
\centering
\resizebox{1\textwidth}{!}{%
\begin{circuitikz}
\tikzstyle{every node}=[font=\LARGE]
\node [font=\LARGE] at (0.75,15.25) {};
\node [font=\LARGE] at (2.5,16.75) {PL};
\draw [short] (2.5,16.25) -- (2.5,15);
\draw  (2.5,14.5) circle (0.5cm);
\draw [short] (2,14.5) -- (3,14.5);
\draw [short] (2.5,15) -- (2.5,14);
\draw [->, >=Stealth] (2.5,12.5) -- (3.75,12.5);
\draw  (3.75,13) rectangle  node {\LARGE F} (5,12);
\draw [->, >=Stealth] (5,12.5) -- (6.25,12.5);
\draw  (6.75,12.5) circle (0.5cm);
\draw [short] (6.25,12.5) -- (7.25,12.5);
\draw [short] (6.75,13) -- (6.75,12);
\draw [->, >=Stealth] (7.25,12.5) -- (8.5,12.5);
\draw  (9,12.5) circle (0.5cm);
\draw [short] (8.5,12.5) -- (9.5,12.5);
\draw [short] (9,13) -- (9,12);
\draw [->, >=Stealth] (9.5,12.5) -- (11,12.5);
\draw  (11,13) rectangle  node {\LARGE L} (12.25,12);
\draw [->, >=Stealth] (12.25,12.5) -- (13.75,12.5);
\draw  (13.75,13) rectangle  node {\LARGE F} (15,12);
\draw [->, >=Stealth] (15,12.5) -- (17,12.5);
\draw  (17.5,12.5) circle (0.5cm);
\draw [short] (17,12.5) -- (18,12.5);
\draw [short] (17.5,13) -- (17.5,12);
\draw [short] (17.5,13) -- (17.5,13.25);
\node [font=\LARGE] at (17.5,16.75) {PR};
\draw [short] (17.5,16.25) -- (17.5,15);
\draw  (17.5,14.5) circle (0.5cm);
\draw [short] (17,14.5) -- (18,14.5);
\draw [short] (17.5,15) -- (17.5,14);
\draw [short] (17.5,14) -- (17.5,13.25);
\draw [short] (2.5,14) -- (2.5,12.5);
\draw [short] (2.5,12.5) -- (2.5,11.25);
\draw [short] (17.5,12) -- (17.5,11.25);
\draw [short] (17.5,11.25) -- (2.5,8.75);
\draw [short] (2.5,11.25) -- (17.5,8.75);
\draw [->, >=Stealth] (2.5,8.75) -- (2.5,8);
\draw [->, >=Stealth] (17.5,8.75) -- (17.5,8);
\node [font=\LARGE] at (1.5,14.5) {KL};
\node [font=\LARGE] at (18.5,14.5) {KR};
\node [font=\LARGE] at (6.75,13.5) {KR};
\node [font=\LARGE] at (9,13.5) {Rc};
\node [font=\LARGE] at (9,13.5) {};
\end{circuitikz}
}%
\caption{ITUbee round structure}
\label{fig:itubee}
\end{figure}

\begin{algorithm}
    \caption{Pseudo code for ITUbee encryption}
    \label{alg:ITUbee}
    \begin{algorithmic}[1]
        \STATE $X_1 \gets P_L \oplus K_L$
        \STATE $X_0 \gets P_R \oplus K_R$
        \FOR{$i = 1$ to $20$}
            \IF{$i \in \{1, 3, \dots, 19\}$}
                \STATE $RK \gets K_R$
            \ELSE
                \STATE $RK \gets K_L$
            \ENDIF
            \STATE $X_{i+1} \gets X_{i-1} \oplus F(L(RK \oplus RC_i \oplus F(X_i)))$
            \STATE \text{Round constant is added to the rightmost 16 bits}
        \ENDFOR
        \STATE $C_L \gets X_{20} \oplus K_R$
        \STATE $C_R \gets X_{21} \oplus K_L$
    \end{algorithmic}
\end{algorithm}

\subsubsection{L Function}
The L funtion which has branch number 4 is computed as,\\
$L(a \parallel b \parallel c \parallel d \parallel e) = (e \oplus a \oplus b) \parallel (a \oplus b \oplus c) \parallel (b \oplus c \oplus d) \parallel (c \oplus d \oplus e) \parallel (d \oplus e \oplus a)$



\subsubsection{F Function}
The F function is computed as,\\
\begin{center}
    $F (X) = S(L(S(X)))$
\end{center}

S is the S-box used in AES \cite{fips197aes}. S is the only non-linear structure in the ITUbee algorithm.  40 bit input of S is splitted into 8 bits then enters the s-box.
\begin{center}
    $S(a\parallel b \parallel c \parallel d \parallel e) = s[a] \parallel s[b] \parallel s[c] \parallel s[d] \parallel s[e]$ 
\end{center}

\subsubsection{Key Schedule}

ITUbee has no key schedule. The algorithm splits the 80 bit key into two 40 bit parts $K_L$ and $K_R$ and uses the right part of the key in odd rounds and the left part of the keys in even rounds.



\subsection{Related-key Differential Cryptanalysis}

Related-key differential cryptanalysis is an extension of classical differential cryptanalysis in which the adversary exploits known relationships between different secret keys. While differential cryptanalysis was originally introduced by Eli Biham and Adi Shamir \cite{biham1990differential}, the related-key attack model was later formalized by Kelsey, Schneier, and Wagner \cite{kelsey1996related}. In this model, the attacker is allowed to observe encryptions under multiple keys that are related through a known difference, and analyzes how both data and key differences propagate through the cipher. The main idea is to construct high-probability differential characteristics that incorporate key differences, enabling the extraction of information about the secret key. This approach is particularly effective against block ciphers with weak or simple key schedule algorithms, where key differences propagate in a predictable manner.

For the related-key model, we must explicitly represent the key differences in addition to the data differences. To do this, we introduce binary variables for key activity and add constraints that force at least one key byte to be active when a related-key difference is present. Without these constraints, the model could avoid key activity and would effectively reduce to the ordinary differential model.

This related-key model also contains more intermediate variables than the differential and linear models, because we must represent the extra variables that appear after the key XOR operation in the middle of the round. The modeling of XOR, linear, and key-dependent components follows the MILP approach of Mouha and Preneel \cite{mouha2012differential}. A sketch of the one-round related-key MILP model is shown in Figure \ref{fig:itubeeRelated}.

\begin{figure}[!ht]
\centering
\resizebox{1\textwidth}{!}{%
\begin{circuitikz}
\tikzstyle{every node}=[font=\large]

\node [font=\large] at (2.5,16) {R1};
\draw (2.5,15.5) to[short] (2.5,14.5);
\draw  (2.5,14) circle (0.5cm);
\draw [short] (2.5,13.5) -- (2.5,11.5);
\draw [->, >=Stealth] (2.5,11.5) -- (4,11.5);
\draw  (4,12) rectangle  node {\large SLS} (5.25,11);
\draw [->, >=Stealth] (5.25,11.5) -- (6.5,11.5);
\draw [short] (2.5,14.5) -- (2.5,13.5);
\draw [short] (2,14) -- (3,14);
\draw  (7,11.5) circle (0.5cm);
\draw [->, >=Stealth] (7.5,11.5) -- (8.75,11.5);
\draw  (8.75,12) rectangle  node {\large L} (10,11);
\draw [->, >=Stealth] (10,11.5) -- (11.5,11.5);
\draw  (11.5,12) rectangle  node {\large SLS} (12.75,11);
\draw [->, >=Stealth] (12.75,11.5) -- (14.5,11.5);
\draw  (15,11.5) circle (0.5cm);
\node [font=\large] at (15,16) {R0};
\draw (15,15.5) to[short] (15,14.5);
\draw  (15,14) circle (0.5cm);
\node [font=\large] at (3,11.75) {};
\draw [short] (15,13.5) -- (15,12);
\draw [short] (6.5,11.5) -- (7.5,11.5);
\draw [short] (7,12) -- (7,11);
\draw [short] (14.5,11.5) -- (15.5,11.5);
\draw [short] (15,12) -- (15,11);
\draw [short] (14.5,14) -- (15.5,14);
\draw [short] (15,14.5) -- (15,13.5);
\draw [short] (15,11) -- (15,10.25);
\draw [short] (2.5,11.5) -- (2.5,10.25);
\draw [short] (2.5,10.25) -- (15,8.25);
\draw [short] (15,10.25) -- (2.5,8.25);
\node [font=\large] at (1.5,14) {KL};
\node [font=\large] at (16,14) {KR};
\node [font=\large] at (7,12.5) {KR};
\node [font=\large] at (3.25,12) {R3};
\node [font=\large] at (5.75,12) {R3a};
\node [font=\large] at (8,12) {R3b};
\node [font=\large] at (10.75,12) {R3c};
\node [font=\large] at (13.75,12) {R3d};
\node [font=\large] at (15.5,10.5) {R4};
\node [font=\large] at (2.5,8) {R4};
\node [font=\large] at (15,8) {R3};
\node [font=\large] at (2,12.5) {R3};
\node [font=\large] at (15.5,12.5) {R2};
\end{circuitikz}
}%
\caption{ITUbee MILP Related-Key Differential Cryptanalysis Sketch}
\label{fig:itubeeRelated}
\end{figure}

We found that 8 round of the algorithm has minimum 16 active s-box. ITUbee algorithm uses AES's \cite{dworkin2001advanced} s-box. AES s-box have best input-output difference $2^{-6}$. Therefore best probability for differential characteristic will be $(2^{-6})^{16}= 2^{-96}  $ which is not usable for related key differential attack. Authors of the ITUbee algorithm claims in \cite{karakoc2013itubee} that algorithm is secure against related key differential cryptanalysis after 10 round. In our model we found that 8 round is enough for security against related key differential cryptanalysis.  Best 8 round characteristic found with the model can be seen in Table \ref{tab:RelatedKeyDif}


\begin{table}[ht]
    \centering
    \begin{tabular}{|c|c|c|}
        \hline
        \textbf{Stage} & \textbf{Binary Index} & \textbf{Characteristic} \\
        \hline
        Key  & KL & \{0: 0.0, 1: 1.0, 2: 0.0, 3: 1.0, 4: 0.0\} \\
                     & KR & \{0: 0.0, 1: 0.0, 2: 0.0, 3: 0.0, 4: 0.0\} \\
        \hline
        1. round  & R1 & \{0: 0.0, 1: 1.0, 2: 0.0, 3: 1.0, 4: 0.0\} \\
        input     & R0 & \{0: 1.0, 1: 0.0, 2: 0.0, 3: 0.0, 4: 1.0\} \\
        \hline
        1. round  & R2 & \{0: 1.0, 1: 0.0, 2: 0.0, 3: 0.0, 4: 1.0\} \\
        middle values & R3-3a-3b-3c-3d & \{0: 0.0, 1: 0.0, 2: 0.0, 3: 0.0, 4: 0.0\} \\     
        \hline
        2. round  & R4 & \{0: 1.0, 1: 0.0, 2: 0.0, 3: 0.0, 4: 1.0\} \\
        input     & R3 & \{0: 0.0, 1: 0.0, 2: 0.0, 3: 0.0, 4: 0.0\} \\
        \hline
        2. round  & R4a & \{0: 0.0, 1: 1.0, 2: 0.0, 3: 1.0, 4: 0.0\} \\
        middle values & R4b-4c-4d & \{0: 0.0, 1: 0.0, 2: 0.0, 3: 0.0, 4: 0.0\} \\
                      
        \hline
        3. round  & R5 & \{0: 0.0, 1: 0.0, 2: 0.0, 3: 0.0, 4: 0.0\} \\
        input     & R4 & \{0: 1.0, 1: 0.0, 2: 0.0, 3: 0.0, 4: 1.0\} \\
        \hline
        3. round middle values     & R5a-5b-5c-5d & \{0: 0.0, 1: 0.0, 2: 0.0, 3: 0.0, 4: 0.0\} \\
        \hline
        4. round  & R6 & \{0: 1.0, 1: 0.0, 2: 0.0, 3: 0.0, 4: 1.0\} \\
        input     & R5 & \{0: 0.0, 1: 0.0, 2: 0.0, 3: 0.0, 4: 0.0\} \\
        \hline
        4. round  & R6a & \{0: 0.0, 1: 1.0, 2: 0.0, 3: 1.0, 4: 0.0\} \\
        middle values & R6b-6c-6d & \{0: 0.0, 1: 0.0, 2: 0.0, 3: 0.0, 4: 0.0\} \\
                              \hline
        5. round  & R7 & \{0: 0.0, 1: 0.0, 2: 0.0, 3: 0.0, 4: 0.0\} \\
        input     & R6 & \{0: 1.0, 1: 0.0, 2: 0.0, 3: 0.0, 4: 1.0\} \\
        \hline
        5. round  & R7a & \{0: 0.0, 1: 0.0, 2: 0.0, 3: 0.0, 4: 0.0\} \\
        middle values & R7b-7c-7d & \{0: 0.0, 1: 0.0, 2: 0.0, 3: 0.0, 4: 0.0\} \\
                      
        \hline
        6. round  & R8 & \{0: 1.0, 1: 0.0, 2: 0.0, 3: 0.0, 4: 1.0\} \\
        input     & R7 & \{0: 0.0, 1: 0.0, 2: 0.0, 3: 0.0, 4: 0.0\} \\
        \hline
        6. round  & R8a & \{0: 0.0, 1: 1.0, 2: 0.0, 3: 1.0, 4: 0.0\} \\
        middle values & R8b-8c-8d & \{0: 0.0, 1: 0.0, 2: 0.0, 3: 0.0, 4: 0.0\} \\
                     
        \hline
        7. round  & R9 & \{0: 0.0, 1: 0.0, 2: 0.0, 3: 0.0, 4: 0.0\} \\
        input     & R8 & \{0: 1.0, 1: 0.0, 2: 0.0, 3: 0.0, 4: 1.0\} \\
        \hline
        7. round middle values & R9a-9b-9c-9d & \{0: 0.0, 1: 0.0, 2: 0.0, 3: 0.0, 4: 0.0\} \\
        \hline
        8. round  & R10 & \{0: 1.0, 1: 0.0, 2: 0.0, 3: 0.0, 4: 1.0\} \\
        input      & R9  & \{0: 0.0, 1: 0.0, 2: 0.0, 3: 0.0, 4: 0.0\} \\
        \hline
        8. round  & R10a & \{0: 0.0, 1: 1.0, 2: 0.0, 3: 1.0, 4: 0.0\} \\
        middle values & R10b-10c-10d & \{0: 0.0, 1: 0.0, 2: 0.0, 3: 0.0, 4: 0.0\} \\
                      
        \hline
        9. round & R11 & \{0: 0.0, 1: 0.0, 2: 0.0, 3: 0.0, 4: 0.0\} \\
        input     & R10 & \{0: 1.0, 1: 0.0, 2: 0.0, 3: 0.0, 4: 1.0\} \\
        \hline
    \end{tabular}
    \caption{Best 8 round related-key differential characteristic of ITUbee algorithm}
    \label{tab:RelatedKeyDif}
\end{table}